\chapter{Опис емпіричних даних}

\section{Загальна характеристика вибірки}

Дослідження охоплює 122 акредитованих освітніх
програм зі спеціальності 022 «Дизайн».

\begin{table}[H]
\centering
\caption{Розподіл програм за рівнем вищої освіти}
\label{tab:programs_by_level}
\begin{tabular}{lcc}
\toprule
Рівень освіти & Кількість програм & Відсоток \\
\midrule
Бакалавр & 74 & 60.7\% \\
Магістр & 43 & 35.2\% \\
Доктор філософії & 5 & 4.1\% \\
\midrule
\textbf{Всього} & \textbf{122} & \textbf{100\%} \\
\bottomrule
\end{tabular}
\end{table}

\section{Структура освітніх компонентів}

Середня кількість освітніх компонентів на програму:
24.9 (SD = 12.9).

\begin{figure}[H]
\centering
\includegraphics[width=0.8\textwidth]{figures/courses_per_program.png}
\caption{Розподіл кількості освітніх компонентів}
\label{fig:courses_dist}
\end{figure}

\section{Наявність силабусів}

Покриття силабусами: 99.0\% компонентів
мають завантажені силабуси.

\section{Типи освітніх компонентів}

\begin{figure}[H]
\centering
\includegraphics[width=0.8\textwidth]{figures/course_types.png}
\caption{Розподіл за типом компонента}
\label{fig:course_types}
\end{figure}
